% Adapted from the existing proper-replay appendix and its archived written proof.
% Written proof accompanied by the Section4.Replay Lean development.
\subsection{Proper generation under replay}
\label{app:replay}
\newcommand{\replayset}[1]{\{#1\}}

This subsection proves the three-language impossibility discussed in
Section~\ref{sec:connections} and the finite-family characterization in
Theorem~\ref{thm:proper-replay}. We first specify the replay model, then
explain why the three-language construction prevents successful generation,
and finally identify exactly when a finite family admits a successful
learner. The model follows~\citet[][P22]{replay}; the three-language family
comes from~\citet[][P31]{memory}. The Lean development verifies the transferred
impossibility, both directions of the characterization, and the finite-state
consequences proved below.

\subsubsection{The replay model}

Let $\mathcal H=(L_1,\ldots,L_N)$, with $N\ge1$, be a finite family of
infinite subsets of a countable universe $U$, and let
$X=\bigcup_{i=1}^N L_i$. Different indices may name the same language.
The target is one fixed member $L_i$ of the family. At round $n\ge0$,
the learner first receives a point $x_n\in X$ and then outputs an index
$j_n\in[N]$. Thus a proper learner always proposes one of the languages
in $\mathcal H$, although that language need not contain every point
observed so far.

A \emph{complete replay text} for target $L_i$ satisfies
\begin{align}
 x_n&\in L_i\cup\bigcup_{t<n}L_{j_t}
       \qquad(n\ge0),\label{app:replay:eq:legal}\\
 L_i&\subseteq\replayset{x_n:n\ge0}.\label{app:replay:eq:full}
\end{align}
The first condition allows each observation to come either from the target
or from any language previously proposed by the learner. The second
condition requires every target point eventually to appear. Observations
may repeat, and there is no bound on how long the adversary may delay a
particular target point. Because there is no earlier output at round zero,
the first observation satisfies $x_0\in L_i$.

The learner succeeds if $L_{j_n}\subseteq L_i$ for all sufficiently large
$n$. It need not identify the target's index or output the whole target
language. We initially allow any deterministic learner that uses its full
input history. No computability restriction is imposed on the family or
on the learner. The finite-state construction below will use only the
membership profile of each observation, namely the set of languages that
contain it.

\subsubsection{The three-language impossibility}
\label{app:replay:construction}

The construction from the bounded-memory paper gives a useful first
example of what replay can hide. After the first output, the same complete
input stream can be legal for two targets that have no common allowable
output. The argument applies to learners with full history.

\begin{suppprop}[Three-language impossibility]
\label{app:replay:prop:triangle}
Let $T\subseteq U$ be infinite, and let $a_1,a_2,a_3\in U\setminus T$
be distinct. The family $L_i=T\cup\replayset{a_s:s\ne i}$, for
$i\in\{1,2,3\}$, does not admit deterministic proper generation in the
limit with replay.
\end{suppprop}

\begin{proof}
Suppose that a learner succeeds for this family. Give it a first point
$x_0\in T$, and let its first output be $j$. Write $p$ and $q$ for the
two indices different from $j$. Continue the input by enumerating every
point of $S=T\cup\replayset{a_1,a_2,a_3}$.

This is a complete legal replay text for either fixed target $L_p$ or
fixed target $L_q$. Indeed, the first point belongs to both targets,
and $S=L_p\cup L_j=L_q\cup L_j$. Every subsequent point is therefore
either a target point or a point of the first output language $L_j$,
which is already eligible for replay. The continuation also contains
every point of each target. Legality uses only the first output; it
places no restriction on the learner's later choices.

The learner sees the same input history in the two cases, so it produces
the same output sequence. If it succeeded for both targets, then after
the larger of the two success times every output language would be
contained in $L_p\cap L_q=T\cup\replayset{a_j}$. But every language in
the family contains two of the three exceptional points, so none is
contained in this intersection. This is a contradiction.
\end{proof}

\subsubsection{The finite-family characterization}

The preceding proof identifies the relevant ambiguity: two targets can
have the same union with the learner's first output language. Equivalently,
they agree outside that language. We now formulate a condition requiring
all targets that remain indistinguishable in this way to admit a common
allowable output.

For $x\in X$, define its membership profile by
$\operatorname{prof}(x)=\replayset{i\in[N]:x\in L_i}$, and let
$\mathcal P=\replayset{\operatorname{prof}(x):x\in X}$.
Every $P\in\mathcal P$ is nonempty. If the first observation has profile
$P$, precisely the indices in $P$ are possible targets at that point.
For a proposed first output $j$, write
$R_j(i)=L_i\setminus L_j$ for the part of target $L_i$ outside that
output. For each $i\in P$, define
\[
 E_j(i;P)=\replayset{k\in P:R_j(k)=R_j(i)}.
\]
This is the group of possible targets having the same residual as $L_i$.
The groups partition $P$. A common allowable output for such a group
is an index $h\in[N]$ satisfying $L_h\subseteq L_k$ for every
$k\in E_j(i;P)$. The language $L_h$ need not equal their intersection,
and its index need not belong to the group.

\begin{suppthm}[Exact finite characterization]
\label{app:replay:thm:main}
The family $\mathcal H$ admits deterministic proper generation in the
limit with replay if and only if the following condition holds:
for every $P\in\mathcal P$, there is an index $j\in[N]$ such that,
for every $i\in P$, there is an index $h\in[N]$ with
\[
 L_h\subseteq\bigcap_{k\in E_j(i;P)}L_k.
\]
The index $j$ need not belong to $P$.
\end{suppthm}

\begin{proof}
\emph{Necessity.}
Fix a realized profile $P$ and choose a first point $x$ with this profile.
Let $j$ be the first output of a learner that succeeds for every target
and every complete legal replay text. Consider any group
$E=E_j(i;P)$, and let $R$ denote the residual shared by its members.
After $x$, supply an enumeration of $S=L_j\cup R$.

For each $k\in E$, we have $x\in L_k$ and
$S=L_j\cup L_k$. Consequently the continuation is legal for target
$L_k$, because the first output permits replay from $L_j$, and it is
complete because it enumerates all of $L_k$. These are different
possible fixed targets for one and the same input sequence. The
learner's outputs are therefore identical in all cases.

By success on target $L_k$, all sufficiently late outputs must be
contained in $L_k$. There are only finitely many indices in $E$, so
we can take the largest of their success times. Every output after
that time is a member of $\mathcal H$ contained in every $L_k$ with
$k\in E$. At least one such output exists, giving the required index
$h$. Since this argument applies to every group, the first output
$j$ satisfies the condition for $P$.

\emph{Sufficiency.}
Assume the stated condition. For each realized profile $P$, fix an
index $j(P)$ satisfying it. Also, for every nonempty set of indices
$E\subseteq[N]$ for which a common allowable output exists, fix one
such index $h(E)$. These choices are made once for the whole family.
In particular, $h(E)$ depends only on $E$, not on the history by which
the learner reaches that set.

On receiving its first observation, the learner sets
$P=\operatorname{prof}(x_0)$, outputs $j=j(P)$, and initializes the
candidate set $C=P$. It keeps the index $j$ for the rest of the run.
For every subsequent observation $x$, it proceeds as follows:
\begin{enumerate}
\item If $x\notin L_j$, replace $C$ by
      $\replayset{k\in C:x\in L_k}$. If $x\in L_j$, leave $C$
      unchanged.
\item Examine the residuals $R_j(k)$ for $k\in C$. If one of them
      is contained in every other surviving residual, let $R$ be
      this least residual and let
      $E=\replayset{k\in C:R_j(k)=R}$. Output $h(E)$.
\item If there is no least residual, output $j$.
\end{enumerate}
On an empty candidate set the learner may output $j$; we will show
that this case never occurs on a legal replay text.

The distinction between a \emph{least} residual and an
\emph{inclusion-minimal} residual is essential here. A least residual
is contained in every surviving residual. An inclusion-minimal one
merely has no strictly smaller survivor; other survivors may be
incomparable with it. The proof needs containment in the true
residual, so the rule uses a least residual and otherwise returns
to $j$.

We first justify that $h(E)$ in the rule is defined. Initially $C=P$.
An update tests membership only at a point outside $L_j$. Two indices
with equal residuals give the same answer to every such test. Thus
an update either retains an entire original group $E_j(i;P)$ or
removes the entire group. By induction, $C$ is always a union of
these original groups. If a least surviving residual exists, its
indices form one entire original group, for which the assumed
condition provides a common allowable output.

We next prove the property that makes the updates reliable. Fix the
true target $i_*$. At every stage, the candidate set contains $i_*$,
and every output language lies inside $L_j\cup L_{i_*}$. Both
assertions hold initially: $x_0\in L_{i_*}$ gives $i_*\in P$, and
the first output is $j$. Suppose they hold for all earlier outputs,
and consider the next observation $x$. By replay legality, $x$
belongs either to $L_{i_*}$ or to an earlier output language. Every
such earlier language lies in $L_j\cup L_{i_*}$ by the induction
hypothesis. Therefore, if $x\notin L_j$, necessarily
$x\in L_{i_*}$. The update retains the true target, so $C$ stays
nonempty.

It remains to check the new output. Outputting $j$ preserves the
claimed containment immediately. Otherwise the learner outputs
$h(E)$ for a least residual $R$. Choose $k\in E$. Since
$L_{h(E)}\subseteq L_k$ and the true target is still in $C$, we have
\[
 L_{h(E)}\setminus L_j
 \subseteq L_k\setminus L_j
 =R
 \subseteq L_{i_*}\setminus L_j.
\]
Hence $L_{h(E)}\subseteq L_j\cup L_{i_*}$, as required. This
completes the induction. In particular, the construction satisfies
\begin{equation}
 L_{j_n}\setminus L_{i_*}\subseteq L_j
       \qquad(n\ge0).\label{app:replay:eq:contain}
\end{equation}
Thus any point outside the target that an output makes available for
replay was already available from the first output. The proof
establishes that later observations outside $L_j$ are trustworthy;
it does not assume this when defining the candidate updates.

Finally, we prove convergence. For every initial candidate $k\in P$
such that $R_j(i_*)\nsubseteq R_j(k)$, choose a point
$w_k\in R_j(i_*)\setminus R_j(k)$. This point belongs to the true
target and lies outside $L_j$, but does not belong to $L_k$.
Completeness guarantees that it eventually appears, at which point
$k$ is removed if it has not been removed already. There are only
finitely many initial candidates, so there is a time after all these
witnesses have appeared. From then on, every surviving residual
contains $R_j(i_*)$. Since $i_*$ survives, its residual is a least
surviving residual.

Moreover, every initial candidate with residual equal to
$R_j(i_*)$ survives forever: each observation used for an update is
a true-target point outside $L_j$, and hence belongs to every one
of these equal residuals. The group selected by the learner from
that time onward is therefore the fixed set $E_j(i_*;P)$. Its
chosen output $h(E_j(i_*;P))$ is contained in $L_{i_*}$, because
$i_*$ is a member of the group. The learner stabilizes to this
correct index, proving sufficiency.
\end{proof}

The proof also gives quantitative information about the learner and
shows that outputs can be restricted to inclusion-minimal languages
in the family. We state these consequences separately to distinguish
them from the existence criterion.

\begin{suppcor}[Finite memory and output choices]
\label{app:replay:cor:implementation}
Whenever the condition of Theorem~\ref{app:replay:thm:main} holds,
there is a learner that reads only membership profiles, uses at
most $1+N2^N$ states, changes its output index at most $N$ times,
and stabilizes to a correct index. It can be chosen so that every
output language is inclusion-minimal in $\mathcal H$, meaning that
no language in $\mathcal H$ is its proper subset. The containment
property~\eqref{app:replay:eq:contain} continues to hold.
\end{suppcor}

\begin{proof}
For the state bound, there is one initial state before the first
observation. Thereafter the learner need only store $(j,C)$, where
$j\in[N]$ and $C\subseteq[N]$. The update tests $j\in
\operatorname{prof}(x)$ and, when this fails, replaces $C$ by
$C\cap\operatorname{prof}(x)$. All residual comparisons and the
choices $j(P)$ and $h(E)$ are fixed data of the finite family.
Thus a finite-state transducer can implement the rule with at most
$1+N2^N$ states. This is an existence statement for the fixed
family, rather than an assertion that an arbitrary representation
allows these data to be computed.

After the first output, the output is a fixed function of $(j,C)$.
Consequently it can change only when $C$ strictly decreases,
apart from a possible change on the first subsequent round when
the initial output $j$ is replaced by the choice for the same
candidate set. Because $C$ starts nonempty and always contains the
true target, it strictly decreases at most $|P|-1\le N-1$ times.
There are therefore at most $N$ output changes. Stabilization and
containment were proved in the theorem.

To obtain inclusion-minimal outputs, first consider any index $j$
that satisfies the theorem's condition for a profile $P$. Since
the family is finite, there is an inclusion-minimal member $L_{j'}$
of the family with $L_{j'}\subseteq L_j$. If two targets agree
outside $L_{j'}$, they also agree outside $L_j$, because
$U\setminus L_j\subseteq U\setminus L_{j'}$. Each group formed
using $j'$ is therefore contained in a group formed using $j$.
A language contained in every target of the larger group is also
contained in every target of the smaller one. Thus $j'$ satisfies
the same condition, and every first-output choice can be made
inclusion-minimal.

Likewise, whenever a group $E$ admits a common allowable language,
choose an inclusion-minimal member of the family below that
language. It is still contained in every target in $E$. Fix these
minimal choices globally as the selectors $h(E)$ and run the same
construction. Every output is now inclusion-minimal, while the
state bound, output-change bound, convergence proof, and containment
argument remain valid.
\end{proof}
